% Exam practice, handout chapters 1-3: moderate and hard questions only.
% Questions: compile as is.  Worked answers: compile with \def\KEY{} (see build.sh).
\documentclass[11pt]{article}
\usepackage{chemguide}
\ifkey\renewcommand{\doctype}{Exam practice ch.~1--3 \textperiodcentered{} answers}
\else\renewcommand{\doctype}{Exam practice ch.~1--3}\fi

% difficulty badges and question heading: \xq{M|H}{title}
\newcommand{\moderate}{\tcbox[on line, colback=BurntOrange!15, colframe=BurntOrange, boxrule=0.4pt,
  arc=0.8mm, left=1pt, right=1pt, top=0pt, bottom=0pt, fontupper=\sffamily\scriptsize]{Moderate}}
\newcommand{\hard}{\tcbox[on line, colback=Maroon!12, colframe=Maroon, boxrule=0.4pt, arc=0.8mm,
  left=1pt, right=1pt, top=0pt, bottom=0pt, fontupper=\sffamily\scriptsize]{Hard \,$\bigstar$}}
\newcounter{xq}
\newcommand{\xpart}{A}
\newcommand{\xq}[2]{\par\needspace{6\baselineskip}\medskip\stepcounter{xq}\noindent
  {\sffamily\bfseries\color{heading}\xpart\thexq.}\enspace
  \if#1M\moderate\else\hard\fi\enspace{\sffamily\bfseries #2}\par\nopagebreak\smallskip}
\newcommand{\xparthead}[2]{\setcounter{xq}{0}\renewcommand{\xpart}{#1}%
  \partheading{Part #1 \,$\cdot$\, #2}}

\begin{document}
% ------------------------------------------------------------------ cover
\thispagestyle{empty}
\vspace*{8mm}
\begin{center}
  {\sffamily\small ENGINEERING \& DATA SCIENCES \,$\cdot$\, CHEMISTRY S1-1}\\[8pt]
  {\sffamily\bfseries\Huge\color{heading} Exam practice}\\[4pt]
  {\sffamily\bfseries\LARGE\color{heading} The atom and its quantum model}\\[8pt]
  {\sffamily\Large Handout chapters 1--3 \ifkey\,$\cdot$\, worked answers\fi}\\[4pt]
  {\normalsize 25 questions \,$\cdot$\, \moderate\ and \hard\ only}
\end{center}
\vspace{5mm}
{\sffamily\bfseries\color{heading} How to use it}
\begin{itemize}[leftmargin=7mm, itemsep=2pt]
\ifkey
  \item Every question is followed by a worked answer in a \textcolor{keyblue}{blue box}. Check the
        method (formula in letters first, then numbers with units) as well as the result.
\else
  \item The questions cover every type of exercise that chapters 1--3 allow: calculations,
        identifications, justifications, data analysis and multi-step problems. Easy recall
        questions are left out; the worksheets cover them.
  \item Work on paper with a calculator, then check with
        \emph{Chem-S1-1\_Exam-Practice-Ch1-3\_Answers.pdf}.
\fi
  \item Moderate: one or two steps with a formula from the handout. Hard: several steps,
        reasoning in reverse, or ideas from two chapters at once.
\end{itemize}

\data{$c = \SI{3.00e8}{\metre\per\second}$\quad $h = \SI{6.62e-34}{\joule\second}$\quad
$e = \SI{1.6e-19}{\coulomb}$\quad $1\ \mathrm{eV} = \SI{1.6e-19}{\joule}$\quad
$N_A = \SI{6.022e23}{\per\mole}$\quad $R_H = \SI{3.29e15}{\hertz}$\quad
$a_0 = 52.9$~pm\quad $k = 1/(4\pi\varepsilon_0) = \SI{8.99e9}{\newton\metre\squared\per\coulomb\squared}$\\
$m_e = \SI{9.1094e-31}{\kilo\gram}$\quad $m_p = 1.0073$~amu\quad $m_n = 1.0087$~amu\quad
$1~\mathrm{amu} = \SI{1.6605e-27}{\kilo\gram}$\quad $g = \SI{9.81}{\metre\per\second\squared}$\quad
$E_n = -13.6\,Z^2/n^2$~eV.}

\vspace{2mm}
{\sffamily\bfseries\color{heading} Map of the question types}\par\vspace{1mm}
{\small
\renewcommand{\arraystretch}{1.12}
\begin{tabular}{@{}p{0.11\linewidth}p{0.62\linewidth}p{0.2\linewidth}@{}}
\toprule
Chapter & Type of question & Questions\\
\midrule
1 & Counting protons, neutrons, electrons; isotopes and isoelectronic species & A1\\
1 & Mass defect and energy of a nucleus & A2\\
1 & Millikan: finding $e$ from data; balancing a charged drop & A3\\
1 & Rutherford: closest approach, density and size of the nucleus & A4\\
1 & Thomson: comparing charge-to-mass ratios & A5\\
2 & Converting between $\lambda$, $\nu$, wavenumber, eV and kJ\,mol$^{-1}$ & B1, B7\\
2 & Identifying a transition from a wavelength & B2\\
2 & Which photons are absorbed; ionisation and the ejected electron & B3, D3\\
2 & Emission cascade: number of lines, extreme wavelengths, visible lines & B4\\
2 & Data analysis: Rydberg's constant from measured lines & B5\\
2 & Hydrogen-like ions; coinciding lines & B6, D2\\
2 & Explaining spectra (lines vs bands, emission vs absorption, colour) & B8\\
3 & Allowed sets of quantum numbers; naming orbitals & C1\\
3 & Counting orbitals and electrons with conditions & C2\\
3 & Nodes; recognising an orbital from its nodes and planes & C3, C8\\
3 & Radial probability density: maximum, probability inside a sphere & C4, C5\\
3 & Degeneracy in H vs many-electron atoms; Madelung order & C6\\
3 & Spin (sodium doublet) and Zeeman splitting & C7\\
3 & Sizes $r = n^2 a_0/Z$ & C9, D2\\
1--3 & Synthesis: Coulomb energy and Bohr levels; unknown ion; photoionisation & D1--D3\\
\bottomrule
\end{tabular}}

% ================================================================== PART A
\clearpage
\xparthead{A}{The atom and its particles (chapter 1)}

\xq{M}{Particle bookkeeping}
(a) Complete the table.\\[2pt]
{\small\begin{tabular}{lccccc}
\toprule
Species & \nuc{37}{17}{Cl}$^-$ & \nuc{56}{26}{Fe}$^{3+}$ & \nuc{18}{8}{O}$^{2-}$ & \nuc{16}{8}{O} &
  \nuc{23}{11}{Na}$^{+}$\\
\midrule
protons / neutrons / electrons & & & & & \\
\bottomrule
\end{tabular}}\\[3pt]
(b) Which species are isotopes of each other? Which are isoelectronic, and with which noble gas?
(c) A species has 26 protons, 30 neutrons and 23 electrons. Write its symbol in full.
\begin{solution}
(a) \nuc{37}{17}{Cl}$^-$: 17 / 20 / 18.\quad \nuc{56}{26}{Fe}$^{3+}$: 26 / 30 / 23.\quad
\nuc{18}{8}{O}$^{2-}$: 8 / 10 / 10.\quad \nuc{16}{8}{O}: 8 / 8 / 8.\quad \nuc{23}{11}{Na}$^+$: 11 / 12 / 10.
(Neutrons $N = A - Z$; electrons $= Z - $ charge.)\\
(b) Isotopes: \iso{18}{O} and \iso{16}{O} (same $Z = 8$, different $A$); the charge does not matter.
Isoelectronic: \nuc{18}{8}{O}$^{2-}$ and \nuc{23}{11}{Na}$^+$, 10 electrons each, like neon.\\
(c) $Z = 26$ is iron, $A = 26 + 30 = 56$, charge $26 - 23 = +3$: \nuc{56}{26}{Fe}$^{3+}$.
\end{solution}

\xq{H}{Where does the mass go? (helium-4)}
A \iso{4}{He} atom has a mass of 4.0026~amu. (a) Compute the total mass of its separate
particles (2 protons, 2 neutrons, 2 electrons; $m_e = 1/1823$~amu). (b) Deduce the mass defect
$\Delta m$ in amu and in kg. (c) The binding energy is $E = \Delta m\,c^2$: compute it per atom, in J
and MeV ($1~\mathrm{MeV} = 1.6\times10^{-13}$~J), and per mole. Compare with a chemical bond
(a few hundred kJ\,mol$^{-1}$).
\begin{solution}
(a) $2\times1.0073 + 2\times1.0087 + 2/1823 = 4.0331$~amu.\\
(b) $\Delta m = 4.0331 - 4.0026 = 0.0305$~amu $= 0.0305\times1.6605\times10^{-27} = 5.06\times10^{-29}$~kg.\\
(c) $E = 5.06\times10^{-29}\times(3.00\times10^8)^2 = 4.56\times10^{-12}$~J $= 28.5$~MeV per atom;
per mole $4.56\times10^{-12}\times6.022\times10^{23} = 2.74\times10^{12}$~J\,mol$^{-1}$, about ten million
times a chemical bond energy: nuclear energies dwarf chemical ones.
\end{solution}

\xq{H}{Millikan's experiment}
(a) Charges measured on oil drops: $6.41$, $9.62$, $3.20$, $12.82$ and $4.81\times10^{-19}$~C. Show that they
are whole multiples of one charge and find it.\\
(b) A drop of oil (density 886~kg\,m$^{-3}$, radius $1.00$~\textmu m) is held still between two
horizontal plates $d = 5.0$~mm apart when $U = 380$~V is applied: the electric force $qU/d$ balances
the weight $mg$. Find its charge and the number of extra electrons it carries.
\begin{solution}
(a) Dividing by the smallest value: $6.41/3.20 = 2.00$, $9.62/3.20 = 3.01$, $12.82/3.20 = 4.01$,
$4.81/3.20 = 1.50$. The ratio 1.50 shows that 3.20 is itself $2e$: dividing everything by
$1.60\times10^{-19}$ gives 4, 6, 2, 8 and 3. All are whole multiples of $e = 1.60\times10^{-19}$~C.\\
(b) $m = \rho\,\tfrac43\pi r^3 = 886\times\tfrac43\pi(1.00\times10^{-6})^3 = 3.71\times10^{-15}$~kg;
$mg = 3.64\times10^{-14}$~N.\quad $q = \dfrac{mg\,d}{U} = \dfrac{3.64\times10^{-14}\times5.0\times10^{-3}}{380}
= 4.79\times10^{-19}$~C $\approx 3e$: three extra electrons.
\end{solution}

\xq{H}{Rutherford on copper}
An $\alpha$ particle (charge $+2e$) with kinetic energy 7.7~MeV heads straight for a copper nucleus
($Z = 29$). (a) At its closest approach all its kinetic energy has become potential energy
$k\,q_1q_2/d$. Find $d$. (b) Why is $d$ only an upper limit for the radius of the nucleus?
(c) The copper nucleus ($A = 63$) has a radius of about 4.8~fm. Compute its density and compare
the radius of the atom (128~pm) with that of its nucleus. (d) Use this ratio to explain why most
$\alpha$ particles cross a metal foil undeflected.
\begin{solution}
(a) $E_k = 7.7\times1.6\times10^{-13} = 1.23\times10^{-12}$~J;\quad
$d = \dfrac{k\,(2e)(29e)}{E_k} = \dfrac{8.99\times10^9\times58\times(1.6\times10^{-19})^2}{1.23\times10^{-12}}
= 1.08\times10^{-14}$~m $= 10.8$~fm.\\
(b) The particle turns back before touching the nucleus; a faster particle would get closer. The
nucleus is at most this large.\\
(c) $m \approx 63\times1.6605\times10^{-27} = 1.05\times10^{-25}$~kg, $V = \tfrac43\pi(4.8\times10^{-15})^3
= 4.6\times10^{-43}$~m$^3$, $\rho \approx 2.3\times10^{17}$~kg\,m$^{-3}$.
$128\times10^{-12}/4.8\times10^{-15} \approx 2.7\times10^4$.\\
(d) The nucleus occupies a fraction $(1/2.7\times10^4)^2 \sim 10^{-9}$ of the cross-section of the
atom: almost every $\alpha$ particle meets only the light, spread-out electrons and goes straight.
\end{solution}

\xq{M}{Thomson's specific charge}
In a deflection tube, the deviation is proportional to $q/m$. Compute $q/m$ for an electron, a proton
and an $\alpha$ particle ($m_\alpha = 4.0015$~amu) and rank their deviations. By what factor is the
electron deviated more than the $\alpha$ particle?
\begin{solution}
Electron: $1.6\times10^{-19}/9.1094\times10^{-31} = 1.76\times10^{11}$~C\,kg$^{-1}$. Proton:
$1.6\times10^{-19}/(1.0073\times1.6605\times10^{-27}) = 9.57\times10^{7}$~C\,kg$^{-1}$. $\alpha$:
$2\times1.6\times10^{-19}/(4.0015\times1.6605\times10^{-27}) = 4.82\times10^{7}$~C\,kg$^{-1}$.
Deviation: electron $\gg$ proton $>$ $\alpha$; the electron is deviated $1.76\times10^{11}/4.82\times10^{7}
\approx 3.6\times10^3$ times more (and in the opposite direction).
\end{solution}

% ================================================================== PART B
\xparthead{B}{Light and the hydrogen spectrum (chapter 2)}

\xq{M}{One photon, five units}
A photon carries 2.48~eV. Give its wavelength, frequency, wavenumber (cm$^{-1}$), the energy of one
mole of such photons (kJ\,mol$^{-1}$), and the region of the spectrum.
\begin{solution}
$E = 2.48\times1.6\times10^{-19} = 3.97\times10^{-19}$~J.\quad $\lambda = hc/E = 5.01\times10^{-7}$~m $= 501$~nm
(visible, blue-green).\quad $\nu = E/h = 5.99\times10^{14}$~Hz.\quad $1/\lambda = 2.00\times10^{6}$~m$^{-1}$
$= 2.00\times10^4$~cm$^{-1}$.\quad Per mole: $3.97\times10^{-19}\times6.022\times10^{23} = 2.39\times10^5$~J\,mol$^{-1}$
$= 239$~kJ\,mol$^{-1}$, the order of a chemical bond.
\end{solution}

\xq{M}{Which transition?}
Hydrogen emits a line at $\lambda = 1282$~nm. Using Rydberg's formula, show that it cannot belong to
the Lyman or the Balmer series, and identify the transition.
\begin{solution}
$\dfrac{1}{n_1^2} - \dfrac{1}{n_2^2} = \dfrac{\nu}{R_H} = \dfrac{c}{\lambda R_H} = \dfrac{3.00\times10^8}{1282\times10^{-9}\times3.29\times10^{15}} = 0.0711$.\\
Lyman ($n_1 = 1$): $1/n_2^2 = 0.929$, $n_2 = 1.04$, not an integer. Balmer ($n_1 = 2$): $1/n_2^2 = 0.179$,
$n_2 = 2.36$, no. Paschen ($n_1 = 3$): $1/n_2^2 = 0.111 - 0.0711 = 0.0400$, $n_2 = 5.00$. The line is
$5 \to 3$ (infrared, as expected for $\lambda > 821$~nm).
\end{solution}

\xq{M}{Photons hitting ground-state hydrogen}
Hydrogen atoms in their ground state receive photons of (a) 12.75~eV, (b) 12.50~eV, (c) 13.22~eV,
(d) 20.0~eV. For each, say whether it is absorbed and what happens. For (d), compute the speed of
the electron that leaves.
\begin{solution}
Absorption needs $h\nu = E_n - E_1 = 13.6(1 - 1/n^2)$ exactly, or $h\nu \ge 13.6$~eV (ionisation).
(a) $1 - 1/n^2 = 0.9375$, $n = 4$: absorbed, $1 \to 4$. (b) $n = 3.52$: not an integer, not absorbed.
(c) $n = 5.98 \approx 6$: absorbed, $1 \to 6$ ($13.6\times35/36 = 13.22$~eV).
(d) $20.0 > 13.6$: ionised, $E_k = 20.0 - 13.6 = 6.4$~eV $= 1.02\times10^{-18}$~J;
$v = \sqrt{2E_k/m_e} = \sqrt{2\times1.02\times10^{-18}/9.1094\times10^{-31}} = 1.50\times10^{6}$~m\,s$^{-1}$.
\end{solution}

\xq{H}{Emission cascade from $n = 5$}
A gas of hydrogen atoms is excited to $n = 5$. (a) How many different lines can it emit? (b) How
many of them are visible, and which? (c) Compute the shortest and the longest wavelengths emitted.
(d) Draw the level diagram with the transitions, grouped into series.
\begin{solution}
(a) Any pair of levels among 1--5: $5\times4/2 = 10$ lines.\\
(b) Visible lines are Balmer lines ending on $n = 2$: $3\to2$ (657~nm), $4\to2$ (486~nm), $5\to2$ (434~nm):
three lines.\\
(c) Shortest: largest energy gap, $5 \to 1$: $\lambda = c/[R_H(1 - 1/25)] = 95.0$~nm (UV). Longest:
smallest gap, $5 \to 4$: $\lambda = c/[R_H(1/16 - 1/25)] = 4.05$~\textmu m (IR).\\
(d) Lyman (to $n = 1$): 4 lines; Balmer (to 2): 3; Paschen (to 3): 2; Brackett (to 4): 1. Total 10.
\end{solution}

\xq{H}{Measuring Rydberg's constant}
The four visible lines of hydrogen are measured at 656.3, 486.1, 434.0 and 410.2~nm.
(a) Show that Balmer's formula predicts a straight line when $1/\lambda$ is plotted against $1/n^2$,
and give its slope and intercept. (b) Using a linear fit (or the first and last points), find $R_H$
in m$^{-1}$ and in Hz. (c) Deduce the ionisation energy of hydrogen in eV and kJ\,mol$^{-1}$, and the
wavelength of the Balmer limit.
\begin{solution}
(a) $\dfrac{1}{\lambda} = \dfrac{R_H}{c}\left(\dfrac14 - \dfrac1{n^2}\right) = \dfrac{R_H}{4c} - \dfrac{R_H}{c}\,\dfrac1{n^2}$:
a line of slope $-R_H/c$ and intercept $R_H/(4c)$, with $n = 3, 4, 5, 6$.\\
(b) Points $(1/n^2, 1/\lambda)$: $(0.1111, 1.5237)$, $(0.0625, 2.0572)$, $(0.0400, 2.3041)$,
$(0.0278, 2.4378)$, with $1/\lambda$ in $10^6$~m$^{-1}$. Least squares: slope $= -1.0972\times10^7$~m$^{-1}$,
intercept $2.743\times10^{6}$~m$^{-1}$ ($\times4 = 1.097\times10^7$, consistent). So
$R_H/c = 1.097\times10^7$~m$^{-1}$ and $R_H = 3.29\times10^{15}$~Hz.\\
(c) $E_I = hR_H = h c\times1.097\times10^7 = 2.18\times10^{-18}$~J $= 13.6$~eV, i.e.\
$2.18\times10^{-18}\times6.022\times10^{23} = 1.31\times10^6$~J\,mol$^{-1} = 1312$~kJ\,mol$^{-1}$.
Balmer limit: $\lambda = 4/(1.097\times10^7) = 365$~nm.
\end{solution}

\xq{H}{The Li$^{2+}$ ion}
(a) Compute the ionisation energy of \ce{Li^2+} in its ground state and the wavelength of its
$2 \to 1$ line. In which region is it? (b) Show that the \ce{Li^2+} transition $9 \to 6$ has
exactly the wavelength of hydrogen's $3 \to 2$ line. (c) Generalise: which \ce{Li^2+} transitions
reproduce all the Balmer lines?
\begin{solution}
(a) $E_I = 13.6\times3^2 = 122.4$~eV. $\nu = 9R_H(1 - 1/4) = 2.22\times10^{16}$~Hz, $\lambda = 13.5$~nm:
far UV (almost X-rays).\\
(b) \ce{Li^2+}: $9\left(\frac1{36} - \frac1{81}\right) = \frac{9}{36} - \frac{9}{81} = \frac14 - \frac19$,
which is hydrogen's $\frac1{2^2} - \frac1{3^2}$: same frequency, same wavelength (657~nm).\\
(c) $9\left(\frac1{n_1^2} - \frac1{n_2^2}\right) = \frac{1}{(n_1/3)^2} - \frac{1}{(n_2/3)^2}$: it equals
$\frac14 - \frac1{m^2}$ for $n_1 = 6$ and $n_2 = 3m$. The transitions $9 \to 6$, $12 \to 6$,
$15 \to 6$\dots\ reproduce the Balmer lines $m = 3, 4, 5\dots$
\end{solution}

\xq{M}{Counting photons}
A 60~W sodium lamp turns 30\,\% of its power into light at 589~nm. How many photons does it emit per
second, and how many moles of photons per hour?
\begin{solution}
Light power $P = 0.30\times60 = 18$~W. One photon: $E = hc/\lambda = 3.37\times10^{-19}$~J.
$N/t = 18/3.37\times10^{-19} = 5.34\times10^{19}$~photons\,s$^{-1}$; per hour
$5.34\times10^{19}\times3600/6.022\times10^{23} = 0.319$~mol of photons.
\end{solution}

\xq{M}{Explaining spectra}
(a) Why do the absorption and emission spectra of hydrogen show lines at exactly the same
wavelengths? (b) Why does a solution of \ce{KMnO4} show a broad absorption band around 530~nm
rather than thin lines, and why does it look purple? (c) What experimental fact led Bohr to say
that the energy of the atom is quantised?
\begin{solution}
(a) Both involve the same pairs of levels: a photon of energy $E_{n_2} - E_{n_1}$ is emitted when the
electron drops and absorbed when it rises.\\
(b) In a molecule in solution, many closely spaced levels (vibrations, interactions with the
solvent) merge into a continuous band. The solution absorbs green light (about 530~nm), so we see
the complementary colour, purple/magenta.\\
(c) Hydrogen absorbs and emits only a few discrete wavelengths: the atom can only exchange definite
amounts of energy, so its energy takes only definite values.
\end{solution}

% ================================================================== PART C
\xparthead{C}{Quantum numbers and atomic orbitals (chapter 3)}

\xq{M}{Allowed or not?}
For each set $(n, \ell, \ml, \ms)$, say whether it is allowed. If so, name the orbital; if not, say
which rule is broken: (a) $(2, 2, 0, +\tfrac12)$; (b) $(3, 2, -2, -\tfrac12)$; (c) $(4, 0, 1, +\tfrac12)$;
(d) $(5, 3, -3, +\tfrac12)$; (e) $(1, 0, 0, 0)$; (f) $(3, 1, -1, +\tfrac12)$.
\begin{solution}
(a) No: $\ell \le n - 1 = 1$. (b) Yes: 3d. (c) No: $|\ml| \le \ell = 0$. (d) Yes: 5f. (e) No: $\ms = \pm\frac12$
only. (f) Yes: 3p.
\end{solution}

\xq{M}{Counting with conditions}
(a) For $n = 4$: list the subshells and give the numbers of orbitals and electrons. (b) How many
electrons of an atom can have $n = 3$ and $\ml = +1$? (c) How many can have $n = 4$ and
$\ms = +\frac12$? (d) How many orbitals with $\ell = 2$ exist for $n \le 5$?
\begin{solution}
(a) 4s, 4p, 4d, 4f: $1 + 3 + 5 + 7 = 16 = n^2$ orbitals, $2n^2 = 32$ electrons.\\
(b) $\ml = +1$ exists for 3p and 3d: 2 orbitals, 4 electrons.\\
(c) 16 orbitals, one electron with $\ms = +\frac12$ each: 16.\\
(d) $\ell = 2$ needs $n \ge 3$: 3d, 4d, 5d, $3\times5 = 15$ orbitals.
\end{solution}

\xq{M}{Nodes}
(a) Give the numbers of radial, angular and total nodes of 2s, 3p, 4d, 5f and 4s. (b) Which orbital
has 2 radial nodes and 1 angular node? Which has no radial node and 3 angular nodes? (c) Can an
orbital of the shell $n = 4$ have 2 radial and 2 angular nodes?
\begin{solution}
(a) Radial $n - \ell - 1$, angular $\ell$, total $n - 1$:
2s 1/0/1;\ 3p 1/1/2;\ 4d 1/2/3;\ 5f 1/3/4;\ 4s 3/0/3.\\
(b) $\ell = 1$, $n = 2 + 1 + 1 = 4$: 4p. $\ell = 3$, $n = 4$: 4f.\\
(c) No: the total would be 4, but a shell $n = 4$ has $n - 1 = 3$ nodes. It is a 5d orbital.
\end{solution}

\xq{H}{Where is the 1s electron?}
For hydrogen, $P(r) = \dfrac{4}{a_0^3}\,r^2\,e^{-2r/a_0}$. (a) Show that $P$ is maximum at $r = a_0$.
(b) Compute $P(2a_0)/P(a_0)$. (c) Show that the probability of finding the electron inside the sphere
of radius $a_0$ is $1 - 5e^{-2} \approx 32\,\%$
(use $\int x^2 e^{-x}\,\mathrm dx = -(x^2 + 2x + 2)\,e^{-x}$). (d) Where is the maximum for \ce{He+}?
\begin{solution}
(a) $P'(r) = \frac{4}{a_0^3}e^{-2r/a_0}\left(2r - \frac{2r^2}{a_0}\right) = 0$ at $r = a_0$ (positive before,
negative after): a maximum.\\
(b) $\dfrac{P(2a_0)}{P(a_0)} = \dfrac{4a_0^2e^{-4}}{a_0^2e^{-2}} = 4e^{-2} = 0.54$.\\
(c) With $x = 2r/a_0$: $\int_0^{a_0}P\,\mathrm dr = \frac12\int_0^{2}x^2e^{-x}\,\mathrm dx
= \frac12\left[2 - 10e^{-2}\right] = 1 - 5e^{-2} = 0.323$.\\
(d) Replace $a_0$ by $a_0/Z$: maximum at $a_0/2 = 26.5$~pm.
\end{solution}

\xq{H}{The 2s orbital and its node}
For the hydrogen 2s orbital, $\Psi_{2s} \propto \left(2 - \dfrac{r}{a_0}\right)e^{-r/2a_0}$.
(a) Find the radius of the radial node, in pm. (b) The radial density is
$P(r) \propto r^2\left(2 - \frac{r}{a_0}\right)^2 e^{-r/a_0}$: show that its maxima are at
$r = (3 \pm \sqrt5)\,a_0$, and say which one is the most probable distance. (c) Where is the node for
\ce{Li^2+}?
\begin{solution}
(a) $\Psi_{2s} = 0$ for $r = 2a_0 = 106$~pm (a sphere; $n - \ell - 1 = 1$ node).\\
(b) Write $P = f^2 e^{-r/a_0}$ with $f = r(2 - r/a_0)$:
$P' = e^{-r/a_0} f\left[2f' - f/a_0\right]$. $f = 0$ gives $r = 0$ and the node $r = 2a_0$ (minima).
$2f' = f/a_0$: $4 - 4r/a_0 = 2r/a_0 - r^2/a_0^2$, i.e.\ $r^2 - 6a_0 r + 4a_0^2 = 0$,
$r = (3 \pm \sqrt5)\,a_0 = 0.76\,a_0$ and $5.24\,a_0$. The outer maximum is the higher one: the most
probable distance is $5.24\,a_0 \approx 277$~pm, consistent with $r = n^2a_0 = 4a_0$ being the size scale.\\
(c) Replace $a_0$ by $a_0/3$: $2a_0/3 = 35$~pm.
\end{solution}

\xq{M}{Degeneracy and Madelung's order}
(a) In hydrogen, how many orbitals have the energy $E_3$? Give their energy in eV. (b) Why does this
degeneracy disappear in a many-electron atom such as Na? Order 3s, 3p, 3d for Na. (c) Use Madelung's
rule to order 6s, 4f, 5d, 6p, 5f, 7s, and say which pairs have the same $n + \ell$.
\begin{solution}
(a) $n^2 = 9$ orbitals (3s, three 3p, five 3d), all at $E_3 = -13.6/9 = -1.51$~eV.\\
(b) The other electrons screen the nucleus; s electrons penetrate closer to it than p and d
electrons, so they feel a larger charge: $E_{3s} < E_{3p} < E_{3d}$.\\
(c) $n + \ell$: 6s 6, 4f 7, 5d 7, 6p 7, 7s 7, 5f 8. Equal sums are ordered by $n$:
6s $<$ 4f $<$ 5d $<$ 6p $<$ 7s $<$ 5f. Same $n + \ell = 7$: 4f, 5d, 6p, 7s.
\end{solution}

\xq{M}{Spin and magnetic field}
(a) Sodium's D line is a doublet at 589.0 and 589.6~nm. Compute the energy gap between the two upper
levels in meV and the wavenumber gap in cm$^{-1}$. What property of the electron explains it?
(b) Into how many sublevels does a 3d level split in a magnetic field (Zeeman effect)? A 4f level?
Which quantum number is revealed?
\begin{solution}
(a) $\Delta E = hc\left(\frac1{589.0\times10^{-9}} - \frac1{589.6\times10^{-9}}\right) = 3.43\times10^{-22}$~J
$= 2.1$~meV; $\Delta(1/\lambda) = 17.3$~cm$^{-1}$. The spin ($\ms = \pm\frac12$) splits the 3p level into two.\\
(b) $2\ell + 1$ values of $\ml$: 5 for 3d, 7 for 4f. It reveals the magnetic quantum number $\ml$.
\end{solution}

\xq{M}{Recognise the orbital}
Name the orbital: (a) one nodal plane, the $(xOy)$ plane, and no radial node; (b) two angular nodes
and one radial node; (c) spherical, with two radial nodes; (d) $\Psinlm{3}{1}{0}$. For (a) and (d), give
$\ml$ and the occupied/unoccupied status in ground-state hydrogen.
\begin{solution}
(a) $\ell = 1$, $n = 2$, plane $(xOy)$: 2p$_z$ ($\ml = 0$). (b) $\ell = 2$, $n = 4$: 4d. (c) $\ell = 0$, $n = 3$:
3s. (d) $n = 3$, $\ell = 1$, $\ml = 0$: 3p$_z$. Both are unoccupied in ground-state hydrogen (only 1s is).
\end{solution}

\xq{H}{Sizes of hydrogen-like orbitals}
Using $r = n^2a_0/Z$: (a) compare the sizes of H 1s, \ce{Li^2+} 2s and \ce{Li^2+} 3d. (b) Which
hydrogen-like ion has a 2p orbital the same size as the 1s of hydrogen? Give the energy of that 2p
level.
\begin{solution}
(a) H 1s: $a_0 = 52.9$~pm. \ce{Li^2+} 2s: $4a_0/3 = 70.5$~pm. \ce{Li^2+} 3d: $9a_0/3 = 3a_0 = 159$~pm.\\
(b) $n^2/Z = 1$ with $n = 2$: $Z = 4$, \ce{Be^3+}. $E_2 = -13.6\times16/4 = -54.4$~eV.
\end{solution}

% ================================================================== PART D
\xparthead{D}{Synthesis: chapters 1 to 3 together}

\xq{H}{Coulomb meets Bohr}
In the Bohr picture the electron of hydrogen orbits at $r_n = n^2a_0$. Its total energy is
$E = -\dfrac{k e^2}{2r}$ (half its Coulomb potential energy). (a) Compute $E$ for $r = a_0$ in J and eV and
compare with $E_1$. (b) Show that this expression gives $E_n = E_1/n^2$. (c) For \ce{He+} in $n = 2$,
give $r$ and $E$; compare with hydrogen's ground state.
\begin{solution}
(a) $E = -\dfrac{8.99\times10^9\times(1.6\times10^{-19})^2}{2\times52.9\times10^{-12}} = -2.18\times10^{-18}$~J
$= -13.6$~eV $= E_1$.\\
(b) $r_n = n^2a_0$ gives $E_n = -\dfrac{ke^2}{2a_0n^2} = \dfrac{E_1}{n^2}$.\\
(c) Hydrogen-like ion: $r = n^2a_0/Z = 4a_0/2 = 106$~pm; $E = -13.6\times4/4 = -13.6$~eV: the same energy
as H in its ground state, but twice as far out.
\end{solution}

\xq{H}{An unknown ion}
A hydrogen-like ion emits a series ending on $n_1 = 2$ whose limit is at 22.8~nm. Identify the ion,
then give its ground-state ionisation energy and the size of its 1s orbital.
\begin{solution}
Limit: $\nu = Z^2R_H/n_1^2$, so $Z^2 = \dfrac{c\,n_1^2}{\lambda R_H} = \dfrac{3.00\times10^8\times4}{22.8\times10^{-9}\times3.29\times10^{15}}
= 16.0$: $Z = 4$, the ion is \ce{Be^3+}. $E_I = 13.6\times16 = 218$~eV; $r_{1s} = a_0/4 = 13.2$~pm.
\end{solution}

\xq{H}{Photoionisation by a helium lamp}
A helium discharge lamp emits UV light at 58.4~nm. (a) Compute the photon energy in eV. (b) It hits
ground-state hydrogen atoms: what happens, and what is the kinetic energy and speed of the electron
that leaves? (c) Could this lamp ionise \ce{He+}? Justify.
\begin{solution}
(a) $E = hc/\lambda = 3.40\times10^{-18}$~J $= 21.3$~eV.\\
(b) $21.3 > 13.6$~eV: the atom is ionised. $E_k = 21.3 - 13.6 = 7.7$~eV $= 1.22\times10^{-18}$~J;
$v = \sqrt{2E_k/m_e} = 1.64\times10^{6}$~m\,s$^{-1}$.\\
(c) No: \ce{He+} needs $13.6\times4 = 54.4$~eV $> 21.3$~eV.
\end{solution}

\end{document}
